\documentclass[xetex,mathserif,serif]{beamer}

\usepackage{xunicode}
\usepackage{xltxtra}
\usepackage{color}
\usepackage{url}
\usepackage{listings}
\usepackage{fontspec}
\usepackage{geometry}
\usepackage{lastpage}
\usepackage{fancyhdr}
\usepackage{amsmath}
\usepackage{amsthm}
\usepackage{amssymb}
\usepackage{blkarray}
\usepackage{multicol}
\usepackage{relsize}

\definecolor{solarized@base03}{HTML}{002B36}
\definecolor{solarized@base02}{HTML}{073642}
\definecolor{solarized@base01}{HTML}{586e75}
\definecolor{solarized@base00}{HTML}{657b83}
\definecolor{solarized@base0}{HTML}{839496}
\definecolor{solarized@base1}{HTML}{93a1a1}
\definecolor{solarized@base2}{HTML}{EEE8D5}
\definecolor{solarized@base3}{HTML}{FDF6E3}
\definecolor{solarized@yellow}{HTML}{B58900}
\definecolor{solarized@orange}{HTML}{CB4B16}
\definecolor{solarized@red}{HTML}{DC322F}
\definecolor{solarized@magenta}{HTML}{D33682}
\definecolor{solarized@violet}{HTML}{6C71C4}
\definecolor{solarized@blue}{HTML}{268BD2}
\definecolor{solarized@cyan}{HTML}{2AA198}
\definecolor{solarized@green}{HTML}{859900}
\definecolor{yaleblue}{HTML}{0E4C92}

\setbeamertemplate{navigation symbols}{}
% \setbeamerfont{title}{family=\old}
% \setbeamerfont{author}{family=\tfont}%
% \setbeamerfont{frametitle}{family=\oldA}
% \setbeamerfont{date}{family=\dfont}

\setbeamertemplate{itemize items}{--}
\setbeamercolor*{item}{fg=black}

\defaultfontfeatures{Mapping=tex-text}
\hypersetup{pdfstartview={FitH}}

\newcommand{\old}[1]{\fontspec[Alternate=1,Ligatures={Common}]{Hoefler Text}\fontsize{18pt}{30pt}\selectfont #1}%
\newcommand{\oldA}[1]{\fontspec[Alternate=1,Ligatures={Common, Rare}]{Hoefler Text}\fontsize{12pt}{15pt}\selectfont #1}%
\newcommand{\oldB}[1]{\fontspec[Ligatures={Common}]{Didot}\fontsize{12pt}{15pt}\color{solarized@base02}\selectfont #1}%
\newcommand{\tfont}[1]{\fontspec[Alternate=1,Ligatures={Common}]{Hoefler Text}\fontsize{12pt}{20pt}\selectfont #1}%
\newcommand{\dfont}[1]{\fontspec[Ligatures={Common}]{Didot}\fontsize{12pt}{12pt}\selectfont #1}%

\newcommand{\minimize}{\mathop{\mathrm{minimize}}}
\newcommand{\argmin}{\mathop{\mathrm{arg\,min}}}
\newcommand{\argmax}{\mathop{\mathrm{arg\,max}}}
\newcommand{\st}{\mathop{\mathrm{subject\,\,to}}}

\newcommand\independent{\protect\mathpalette{\protect\independenT}{\perp}}
\def\independenT#1#2{\mathrel{\rlap{$#1#2$}\mkern2mu{#1#2}}}

\setlength{\parindent}{0pt}
\setlength{\parskip}{12pt}

\setromanfont [Ligatures={Common}, Numbers={OldStyle}, Variant=01,
 BoldFont={LinLibertine_RB.otf},
 ItalicFont={LinLibertine_RI.otf},
 BoldItalicFont={LinLibertine_RBI.otf}
 ]{LinLibertine_R.otf}



\title{high performance data i/o}
\date{2015-02-16}

\begin{document}

%%%%%%%%%%%%%%%%%%%%%%%%%%%%%%%%%%%%%%%%%%%%%%%%%%%
\begin{frame}[fragile] \frametitle{}

\vfill

{\fontsize{0.7cm}{0cm}\selectfont Lecture 03 \\\vspace{0.2cm} Simple Linear Models:\\
Leverage, Hypothesis Tests, Goodness of Fit}\\\vspace{0.5cm}
09 September 2015

\vspace{2cm}

\begin{minipage}{0.6\textwidth}
Taylor B. Arnold \\
Yale Statistics \\
STAT 312/612
\end{minipage}
\hfill
\begin{minipage}{0.3\textwidth}\raggedleft
\includegraphics[scale=0.3]{../yale-logo.png}
\end{minipage}%

\end{frame}


%%%%%%%%%%%%%%%%%%%%%%%%%%%%%%%%%%%%%%%%%%%%%%%%%%%
\begin{frame}[fragile] \frametitle{}

{\color{yaleblue}\fontsize{16pt}{20pt}\selectfont Notes}

\begin{itemize}
\item Problem Set \#1 Online: Due Next Wednesday, 2015-09-16
\item R code online
\item Course Pace
\end{itemize}

\end{frame}

%%%%%%%%%%%%%%%%%%%%%%%%%%%%%%%%%%%%%%%%%%%%%%%%%%%
\begin{frame}[fragile] \frametitle{}

{\color{yaleblue}\fontsize{16pt}{20pt}\selectfont Goals for today}

\begin{enumerate}
\item simulation of leverage
\item hypothesis tests for simple linear regression
\item goodness of fit, $R^2$
\item Galton's heights data
\end{enumerate}

\end{frame}

%%%%%%%%%%%%%%%%%%%%%%%%%%%%%%%%%%%%%%%%%%%%%%%%%%%
\begin{frame}[fragile] \frametitle{}

\begin{flushright}
{\color{yaleblue}\sc\fontsize{1cm}{0cm}\selectfont Leverage Simulation}
\end{flushright}

\end{frame}

%%%%%%%%%%%%%%%%%%%%%%%%%%%%%%%%%%%%%%%%%%%%%%%%%%%
\begin{frame}[fragile] \frametitle{}

\begin{flushright}
{\color{yaleblue}\sc\fontsize{1cm}{0cm}\selectfont Hypothesis Tests}
\end{flushright}

\end{frame}

%%%%%%%%%%%%%%%%%%%%%%%%%%%%%%%%%%%%%%%%%%%%%%%%%%%
\begin{frame}[fragile] \frametitle{}

{\color{yaleblue}\fontsize{16pt}{20pt}\selectfont Z-Test}

Take the simple linear regression model:
\begin{align*}
y_i &=  x_i\beta  + \epsilon_i, \quad i = 1, \ldots n.
\end{align*}
With independent, identically distributed normal error terms:
\begin{align*}
\epsilon_i \sim \mathcal{N}(0, \sigma^2)
\end{align*}

\end{frame}

%%%%%%%%%%%%%%%%%%%%%%%%%%%%%%%%%%%%%%%%%%%%%%%%%%%
\begin{frame}[fragile] \frametitle{}

Last time we calculated the MLE estimator,
\begin{align*}
\widehat{\beta}_{MLE} &= \frac{\sum_i y_i x_i}{\sum_i x_i^2}
\end{align*}
And showed that it has a normally distribution with the following
mean and variance:
\begin{align*}
\widehat{\beta} \sim \mathcal{N} (\beta, \frac{\sigma^2}{\sum_i x_i^2})
\end{align*}

\end{frame}

%%%%%%%%%%%%%%%%%%%%%%%%%%%%%%%%%%%%%%%%%%%%%%%%%%%
\begin{frame}[fragile] \frametitle{}

If we want to test the hypothesis $H_0: \beta = b$, we
could construct a test statistic as follows: \pause
\begin{align*}
z &= \frac{\widehat{\beta} - b}{\sqrt{\frac{\sigma^2}{\sum_i x_i^2}}}
\end{align*}
\pause Under the null hypothesis, we have
\begin{align*}
z | H_0  \sim \mathcal{N} (0, 1)
\end{align*}

\end{frame}

%%%%%%%%%%%%%%%%%%%%%%%%%%%%%%%%%%%%%%%%%%%%%%%%%%%
\begin{frame}[fragile] \frametitle{}

When we know $\sigma^2$ that is all we need to do, however
outside of simulations we (very) rarely known the true variance
of the noise. Otherwise, we first need to estimate it.

\end{frame}

%%%%%%%%%%%%%%%%%%%%%%%%%%%%%%%%%%%%%%%%%%%%%%%%%%%
\begin{frame}[fragile] \frametitle{}

{\color{yaleblue}\fontsize{16pt}{20pt}\selectfont T-Test}

The residuals from a given prediction of $\beta$ are given by:
\begin{align*}
r_i &= y_i - \widehat{y}_i \\
&= y_i - x_i \widehat{\beta}
\end{align*}
These represent an estimate of the error terms $\epsilon_i$.

\end{frame}

%%%%%%%%%%%%%%%%%%%%%%%%%%%%%%%%%%%%%%%%%%%%%%%%%%%
\begin{frame}[fragile] \frametitle{}

If $r_i$ is the sampled and estimated version of $\epsilon_i$,
it would seem reasonable to have:
\begin{align*}
\frac{1}{n} \sum_{i=1}^n r_i^2 &\approx \mathbb{E} \epsilon^2 \\
&= \sigma^2
\end{align*}

\end{frame}

%%%%%%%%%%%%%%%%%%%%%%%%%%%%%%%%%%%%%%%%%%%%%%%%%%%
\begin{frame}[fragile] \frametitle{}

Much like when estimating the mean a randomly sampled normal
distribution, this is approximately correct though the exact
formula requires a small correction since\pause
\begin{align*}
\mathbb{E} \left( \sum_i r_i^2 \right) &= (n-1) \cdot \sigma^2
\end{align*}
\pause I will delay a formal derivation of this until the multivariate
case; conceptually seems reasonable that the estimate will be slightly
smaller due to the estimation of $r_i$ by the same data.

\end{frame}


%%%%%%%%%%%%%%%%%%%%%%%%%%%%%%%%%%%%%%%%%%%%%%%%%%%
\begin{frame}[fragile] \frametitle{}

So, we instead use a corrected form to estimate the error
variance, an estimator that we will call $s^2$:
\begin{align*}
s^2 &= \frac{1}{n-1} \cdot \sum_i r_i^2 \\
&= \frac{1}{n-1} \cdot \sum_i (y_i - \widehat{y}_i)^2 \\
&= \frac{1}{n-1} \cdot \sum_i (y_i - x_i \beta)^2
\end{align*}

\end{frame}

%%%%%%%%%%%%%%%%%%%%%%%%%%%%%%%%%%%%%%%%%%%%%%%%%%%
\begin{frame}[fragile] \frametitle{}

The ratio of our estimator to the true variance has
a $\chi^2$ distribution with $n-1$ degrees of freedom.
\begin{align*}
\frac{s^2}{\sigma^2} \sim \chi^2_{n - 1}
\end{align*}

\end{frame}

%%%%%%%%%%%%%%%%%%%%%%%%%%%%%%%%%%%%%%%%%%%%%%%%%%%
\begin{frame}[fragile] \frametitle{}

The standard error is then given by:
\begin{align*}
\text{S.E.}(\widehat{\beta}) &= \sqrt{\frac{s^2}{\sum_i x_i^2}} \\
&= \sqrt{\frac{(y - x_i \widehat{\beta})^2}{(n-1) \cdot \sum_i x_i^2}}
\end{align*}

\end{frame}

%%%%%%%%%%%%%%%%%%%%%%%%%%%%%%%%%%%%%%%%%%%%%%%%%%%
\begin{frame}[fragile] \frametitle{}

Finally, we can construct a test statistic:
\begin{align*}
t &= \frac{\widehat{\beta} - b}{\text{S.E.}(\widehat{\beta})}
\end{align*}
\pause And under the null hypothesis, we have
\begin{align*}
t | H_0  \sim t_{n-1}
\end{align*}

\end{frame}

%%%%%%%%%%%%%%%%%%%%%%%%%%%%%%%%%%%%%%%%%%%%%%%%%%%
\begin{frame}[fragile] \frametitle{}

On a related note, we can similarly calculate a confidence
interval for $\beta$ using the standard error. A $100(1-\alpha)\%$.
confidence interval is given by:
\begin{align*}
\widehat{\beta} \pm t_{n-1, 1 - \alpha/2} \cdot \text{S.E.}(\widehat{\beta})
\end{align*}
\pause For a reasonably large sample size $n$, we can approximate
this by a normal distribution:
\begin{align*}
\widehat{\beta} \pm z_{1 - \alpha/2} \cdot \text{S.E.}(\widehat{\beta})
\end{align*}

\end{frame}

%%%%%%%%%%%%%%%%%%%%%%%%%%%%%%%%%%%%%%%%%%%%%%%%%%%
\begin{frame}[fragile] \frametitle{}

{\color{yaleblue}\fontsize{16pt}{20pt}\selectfont F-Test}

As an alternative to the T-test, consider squaring the
test statistic
\begin{align*}
T^2 &= \left( \frac{\widehat{\beta} - b}{\text{S.E.}(\widehat{\beta})} \right)^2 \\
&= \frac{{\color{solarized@red}\left(\frac{\widehat{\beta} - b}{\sigma^2 / \sum_i x_i^2}\right)^2}}{{\color{solarized@violet}s^2 / \sigma^2}} \\
&= \frac{{\color{solarized@red}U}}{{\color{solarized@violet}V}}
\end{align*}
\pause Where ${\color{solarized@red}U} \sim \chi^2_{1}$ and ${\color{solarized@violet}V} \sim \chi^2_{n-1}$.

\pause And therefore $T^2 \sim F_{1, n-1}$.

\end{frame}

%%%%%%%%%%%%%%%%%%%%%%%%%%%%%%%%%%%%%%%%%%%%%%%%%%%
\begin{frame}[fragile] \frametitle{}

{\color{yaleblue}\fontsize{16pt}{20pt}\selectfont Intercept Model}

When we have the model $y = \alpha + x \beta + \epsilon$, the form of $s^2$
changes slightly:
\begin{align*}
s^2 &= \frac{1}{n-2} \cdot \sum_i (y_i - \widehat{y}_i)^2
\end{align*}
\pause as well as the standard errors:
\begin{align*}
\text{S.E.}(\alpha) &= \sqrt{s^2 \cdot \left(\frac{1}{n} + \frac{\bar{x}}{\sum_i (x_i-\bar{x})^2} \right)} \\
\text{S.E.}(\beta) &= \sqrt{\frac{s^2}{\sum_i (x_i-\bar{x})^2}}
\end{align*}

\end{frame}

%%%%%%%%%%%%%%%%%%%%%%%%%%%%%%%%%%%%%%%%%%%%%%%%%%%
\begin{frame}[fragile] \frametitle{}

\begin{flushright}
{\color{yaleblue}\sc\fontsize{1cm}{0cm}\selectfont Goodness of Fit}
\end{flushright}

\end{frame}

%%%%%%%%%%%%%%%%%%%%%%%%%%%%%%%%%%%%%%%%%%%%%%%%%%%
\begin{frame}[fragile] \frametitle{}

{\color{yaleblue}\fontsize{16pt}{20pt}\selectfont $R^2$}

A common measurement of how well a linear model explains
the data is the $R^2$. For the non-intercept version, it
can be written as:
\begin{align*}
R^2 &= 1 - \frac{\sum_i (y_i - \widehat{y}_i)^2}{\sum_i (y_i)^2}
\end{align*}
\pause The more typically seen version compares the estimated
residuals with the centered values of $y$.
\begin{align*}
R^2 &= 1 - \frac{\sum_i (y_i - \widehat{y}_i)^2}{\sum_i (y_i - \bar{y})^2}
\end{align*}

\end{frame}

%%%%%%%%%%%%%%%%%%%%%%%%%%%%%%%%%%%%%%%%%%%%%%%%%%%
\begin{frame}[fragile] \frametitle{}

With a bit of algebraic manipulation, we see that this is
equal to the squared sample correlation of $x$ and $y$:
\begin{align*}
R^2 &= \left( \frac{\sum_i (x_i - \bar{x})(y_i - \bar{y})}{\sqrt{\sum_i (x_i - \bar{x})^2 \cdot \sum_i (y_i - \bar{y})^2}} \right)^2 \\
&= cor(x,y)^2
\end{align*}

\end{frame}

%%%%%%%%%%%%%%%%%%%%%%%%%%%%%%%%%%%%%%%%%%%%%%%%%%%
\begin{frame}[fragile] \frametitle{}

\begin{flushright}
{\color{yaleblue}\sc\fontsize{1cm}{0cm}\selectfont Applications}
\end{flushright}

\end{frame}


\end{document}













